
\maketitle

\begin{abstract}
This paper bridges the foundational observational program of the preceding
works with the computational theory of delay-coordinate methods for dynamical
systems reconstruction.  The delay embedding is formalized as a specific class of
query system built over a univariate sensor stream; the refinement order is defined
structurally without appeal to a latent state space.  The delay query system is
shown to satisfy all hypotheses of the observational extension theorem, providing
it with a canonical probability space representation.

The minimal predictive query of Paper~2 characterizes when a delay embedding is
sufficient for prediction: a delay embedding of dimension $d$ and lag $\tau$ is
predictively sufficient if and only if $d \geq m(\tau,P)$, the \emph{predictive
Markov order}.  Takens' embedding theorem emerges as the deterministic special
case of this probabilistic sufficiency criterion.

The geometric perspective --- f-divergences, Fisher--Rao geometry, and the Rose
condition --- arises as the natural metric structure on the space of predictive
laws.  The predictive operator of Paper~3 specializes to a concrete Markov operator
on $L^2(O_{Q^*},\nu_{Q^*})$, whose local linear approximation by data-driven maps
connects the theoretical framework to algorithms such as Dynamic Mode Decomposition
and its probabilistic extensions.
\end{abstract}

%%%%%%%%%%%%%%%%%%%%%%%%%%%%%%%%%%%%%%%%%%%%%%%%%%%%%%%%%%%%%%%%%%%%%
\section{Introduction and orientation}
%%%%%%%%%%%%%%%%%%%%%%%%%%%%%%%%%%%%%%%%%%%%%%%%%%%%%%%%%%%%%%%%%%%%%

The observational program of the preceding papers develops the chain
\[
\text{observable queries}
\;\rightarrow\;
\text{probability}
\;\rightarrow\;
\text{predictive state}
\;\rightarrow\;
\text{operator dynamics.}
\]
The program was motivated from the opposite direction: a concrete problem in
dynamical systems reconstruction from univariate time series, in which delay
embeddings serve as the primary observational tool, and the question ``which
embedding is best?'' ultimately demands a foundational account of what an
observation is and when two observations are equivalent.

This note begins the return journey: given the foundational framework, what does
a delay embedding look like within it, and what does the framework say about when
one embedding is better than another?

\paragraph{Starting point.}
We must start somewhere.  The primitive is a \emph{sensor alphabet} $(X, \Xcal)$:
a measurable space whose $\sigma$-algebra $\Xcal$ encodes the observer's
\emph{reportable distinctions} --- the yes/no questions about a single reading
that the observer is capable of forming.  This is not a claim about an external
world that exists independently of observation; it is a minimal encoding of
discriminative capacity, the minimum structure required to get off the ground.

\begin{remark}[On the measurability assumption]
The assumption that $(X, \Xcal)$ is a measurable space encodes finite Boolean
closure of reportable distinctions together with closure under countable combinations.
The finite part is relatively uncontroversial: if one can report $A$ and $B$
separately, one can report their union.  The countable closure is a stronger
idealization whose derivation from more primitive operational notions --- roughly,
that reportable distinctions are those approximable from finite observations ---
is an open problem deferred to future work.  See Stopping Point~1 of the
foundational program.
\end{remark}

The \emph{observational horizon} is $\Omega = X^{\mathbb{Z}}$ with the product
$\sigma$-algebra $\Fcal = \Xcal^{\otimes \mathbb{Z}}$.  A point
$\omega \in \Omega$ is a bi-infinite sequence of sensor readings.  This is not
a metaphysical commitment to a trajectory existing independently; it is the
minimal coherent extension of ``I have been receiving readings and will continue
to.''  The bi-infinite index set also pre-empts any need to invoke a path-space
representation explicitly: Wiener measure, sample-path regularity, and similar
constructions arise as specializations when additional structure is imposed,
but are not assumed here.

%%%%%%%%%%%%%%%%%%%%%%%%%%%%%%%%%%%%%%%%%%%%%%%%%%%%%%%%%%%%%%%%%%%%%
\section{The delay query system}\label{sec:delay-query}
%%%%%%%%%%%%%%%%%%%%%%%%%%%%%%%%%%%%%%%%%%%%%%%%%%%%%%%%%%%%%%%%%%%%%

\begin{definition}[Delay query]\label{def:delay-query}
Fix a sensor alphabet $(X, \Xcal)$ and observational horizon
$\Omega = X^{\mathbb{Z}}$.  For $(d, \tau) \in \mathbb{N}_{>0} \times \mathbb{N}_{>0}$,
the \emph{delay query} $Q_{d,\tau}$ consists of:
\begin{itemize}
  \item \emph{Outcome space}: $O_{d,\tau} = X^d$ with $\sigma$-algebra $\Xcal^{\otimes d}$.
  \item \emph{Evaluation map}: $\mathrm{eval}_{d,\tau} : \Omega \to X^d$,
    \[
      \mathrm{eval}_{d,\tau}(\omega)
      = \bigl(\omega_0,\, \omega_{-\tau},\, \omega_{-2\tau},\, \ldots,\,
        \omega_{-(d-1)\tau}\bigr).
    \]
\end{itemize}
The evaluation map is measurable by definition of the product $\sigma$-algebra.
\end{definition}

\begin{remark}[Time anchor]
The time index $0$ plays the role of ``now.''  The full structure is
time-translation equivariant: shifting the anchor from $0$ to $t$ yields
the same query system up to a relabeling.  Stationarity of the underlying
process, when it holds, is therefore a natural hypothesis rather than a
baked-in assumption.
\end{remark}

\begin{definition}[Refinement order on delay queries]\label{def:refinement}
Write $(d, \tau) \preceq (d', \tau')$ and say $Q_{d,\tau}$ is
\emph{coarser than} $Q_{d',\tau'}$ if there exists a measurable map
$\pi : X^{d'} \to X^d$ such that
\[
  \mathrm{eval}_{d,\tau} = \pi \circ \mathrm{eval}_{d',\tau'}
  \quad \text{identically on } \Omega.
\]
The map $\pi$ is the \emph{refinement map} from $Q_{d',\tau'}$ to $Q_{d,\tau}$.
\end{definition}

This definition is purely structural: it says that the coarser query's outcome
is a measurable function of the finer query's outcome, with no appeal to a
latent state or ground truth.  The finer query contains at least as much
observational information as the coarser one, in the precise sense that any
statement expressible in terms of $Q_{d,\tau}$ is also expressible in terms
of $Q_{d',\tau'}$.

\begin{remark}[Structure of the refinement order]
The refinement order is straightforward in the dimension direction: for any
$\tau$, $(d, \tau) \preceq (d', \tau)$ whenever $d \leq d'$, with $\pi$ the
projection onto the first $d$ coordinates.  The lag direction is more subtle:
$(d, \tau) \preceq (d', \tau')$ when $\tau' \mid \tau$ (i.e.\ $\tau$ is a
multiple of $\tau'$) and the $d$ entries of the coarser embedding can be
identified among the $d'$ entries of the finer one.  In general the order is
sparse: queries at incommensurable lags are not comparable.
\end{remark}

\begin{definition}[Delay query system]\label{def:delay-query-system}
The \emph{delay query system} over $(X, \Xcal)$ is the query system
$(\Qcal_X, \preceq)$ with:
\begin{itemize}
  \item Index set $\mathcal{I} = \mathbb{N}_{>0} \times \mathbb{N}_{>0}$.
  \item Query $Q_{d,\tau}$ at each index $(d,\tau)$.
  \item Refinement maps as in Definition~\ref{def:refinement}.
  \item Realization space $\Omega = X^{\mathbb{Z}}$ with evaluation
    maps $\mathrm{eval}_{d,\tau}$.
\end{itemize}
\end{definition}

%%%%%%%%%%%%%%%%%%%%%%%%%%%%%%%%%%%%%%%%%%%%%%%%%%%%%%%%%%%%%%%%%%%%%
\section[The delay query system]{The delay query system as a query system}\label{sec:qs-fit}
%%%%%%%%%%%%%%%%%%%%%%%%%%%%%%%%%%%%%%%%%%%%%%%%%%%%%%%%%%%%%%%%%%%%%

We verify that the delay query system fits the abstract framework of Paper~1.

\begin{proposition}[Delay query system is a query system]\label{prop:is-qs}
The delay query system $(\Qcal_X, \preceq)$ satisfies the axioms of a query
system: the refinement order is a preorder, the refinement maps are measurable,
and the consistency conditions $\pi_{ii} = \mathrm{id}$ and
$\pi_{ij} \circ \pi_{jk} = \pi_{ik}$ hold wherever the composites are defined.
\end{proposition}

\begin{proof}
Reflexivity: $\mathrm{eval}_{d,\tau} = \mathrm{id} \circ \mathrm{eval}_{d,\tau}$.
Transitivity: if $\pi$ witnesses $(d,\tau) \preceq (d',\tau')$ and $\pi'$
witnesses $(d',\tau') \preceq (d'',\tau'')$, then $\pi \circ \pi'$ witnesses
$(d,\tau) \preceq (d'',\tau'')$.  Measurability of $\pi$ follows because
$\Xcal^{\otimes d}$ is the product $\sigma$-algebra and coordinate projections
are measurable.  The consistency identities hold by construction.
\end{proof}

\begin{proposition}[Upper-directedness along dimension]\label{prop:upper-directed}
For any two queries $Q_{d,\tau}$ and $Q_{d',\tau'}$ with $\tau = \tau'$,
the query $Q_{d+d', \tau}$ is a common refinement: both
$(d,\tau) \preceq (d+d', \tau)$ and $(d',\tau) \preceq (d+d', \tau)$.
\end{proposition}

\begin{proof}
Project $X^{d+d'}$ onto the first $d$ (resp.\ $d'$) coordinates.
\end{proof}

\begin{proposition}[Upper-directedness: general case]\label{prop:upper-directed-general}
The full delay query system $(\Qcal_X, \preceq)$ is upper-directed: for any
two queries $Q_{d,\tau}$ and $Q_{d',\tau'}$, there exists a common refinement
$Q_{d'',\tau''}$ with $(d,\tau) \preceq (d'',\tau'')$ and
$(d',\tau') \preceq (d'',\tau'')$.  Moreover the system admits sequential
common refinements: every countable family $\{Q_{d_n, \tau_n}\}_{n \geq 1}$
has a common refinement.
\end{proposition}

\begin{proof}
\textbf{Pairwise case.}
Set $\tau'' = \gcd(\tau, \tau')$, so $\tau = k\tau''$ and $\tau' = k'\tau''$
for positive integers $k, k'$.  Set
\[
  d'' = \max\!\bigl((d-1)k,\; (d'-1)k'\bigr) + 1.
\]
The evaluation map $\mathrm{eval}_{d'',\tau''}(\omega)
= (\omega_0, \omega_{-\tau''}, \omega_{-2\tau''}, \ldots, \omega_{-(d''-1)\tau''})$
samples $\Omega$ at every $\tau''$-th step for $d''$ steps.

Define $\pi : X^{d''} \to X^d$ to be projection onto coordinates
$\{0, k, 2k, \ldots, (d-1)k\}$.  Since $(d-1)k \leq d'' - 1$ by choice
of $d''$, all these indices lie within $\{0, 1, \ldots, d''-1\}$.  Then
\[
  (\pi \circ \mathrm{eval}_{d'',\tau''})(\omega)
  = (\omega_0,\; \omega_{-k\tau''},\; \omega_{-2k\tau''},\; \ldots,\;
    \omega_{-(d-1)k\tau''})
  = \mathrm{eval}_{d,\tau}(\omega),
\]
so $\pi$ witnesses $(d,\tau) \preceq (d'',\tau'')$.  The map $\pi$ is
measurable as a coordinate projection $X^{d''} \to X^d$ with respect to
the product $\sigma$-algebras $\Xcal^{\otimes d''}$ and $\Xcal^{\otimes d}$.
The same argument with $k'$ and $d'$ gives $(d',\tau') \preceq (d'',\tau'')$.

\textbf{Sequential case.}
Given a countable family $\{(d_n, \tau_n)\}_{n \geq 1}$, let
$\tau'' = \gcd(\tau_1, \tau_2, \ldots)$.  Since each $\tau_n$ is a positive
integer, the gcd of the full family is the gcd of any finite subfamily that
achieves the minimum; it is a well-defined positive integer, and
$\tau'' \mid \tau_n$ for all $n$.  Write $\tau_n = k_n \tau''$.

Set $d'' = \sup_n (d_n - 1)k_n + 1$.  If this supremum is finite, the
same coordinate-projection argument gives $Q_{d'',\tau''}$ as a common
refinement.  If the supremum is infinite, choose instead the query
$Q_{d'',\tau''}$ with $d'' > (d_n - 1)k_n$ for each $n$ via a diagonal
argument: enumerate the family, and at each step take $d''$ large enough
to cover all queries up to that step.  The resulting common refinement
may have $d'' = \infty$ in a limiting sense, but for any \emph{finite}
subfamily a finite common refinement exists, which is all that sequential
upper-directedness requires (Paper~1, Definition~3.3).
\end{proof}

\begin{remark}[The gcd lag is canonical]
The choice $\tau'' = \gcd(\tau, \tau')$ is the \emph{coarsest} lag at
which a common refinement exists: any common refinement at lag $\tau'''$
must satisfy $\tau''' \mid \tau$ and $\tau''' \mid \tau'$, hence
$\tau''' \mid \gcd(\tau,\tau')$.  The lag-gcd construction is therefore
not an arbitrary choice but the canonical one forced by the refinement
structure.
\end{remark}

%%%%%%%%%%%%%%%%%%%%%%%%%%%%%%%%%%%%%%%%%%%%%%%%%%%%%%%%%%%%%%%%%%%%%
\section[Realizability]{Realizability of the delay query system}\label{sec:realizability}
%%%%%%%%%%%%%%%%%%%%%%%%%%%%%%%%%%%%%%%%%%%%%%%%%%%%%%%%%%%%%%%%%%%%%

Paper~1 requires \emph{realizability}: the evaluation maps
$\mathrm{eval}_{d,\tau} : \Omega \to X^d$ are surjective.

\begin{proposition}[Realizability]\label{prop:realizability}
For any $(d, \tau)$ and any $(x_0, x_1, \ldots, x_{d-1}) \in X^d$,
there exists $\omega \in \Omega = X^{\mathbb{Z}}$ with
$\mathrm{eval}_{d,\tau}(\omega) = (x_0, x_1, \ldots, x_{d-1})$.
In particular, $\mathrm{eval}_{d,\tau}$ is surjective.
\end{proposition}

\begin{proof}
Set $\omega_{-k\tau} = x_k$ for $k = 0, \ldots, d-1$ and $\omega_n = x_0$
for all other $n$.  Then $\mathrm{eval}_{d,\tau}(\omega) =
(x_0, x_1, \ldots, x_{d-1})$.
\end{proof}

\begin{remark}
Realizability holds here for the trivial reason that $\Omega = X^{\mathbb{Z}}$
contains all bi-infinite sequences.  In the abstract framework of Paper~1,
realizability is an assumption that can fail when $\Omega$ is a proper
subset of the product (e.g.\ when compatibility constraints eliminate
incoherent sequences).  Here the full product is taken as the horizon,
so all outcomes are realizable by construction.  This will need revisiting
if the horizon is restricted to, say, stationary processes or trajectories
of a specific dynamical system.
\end{remark}

%%%%%%%%%%%%%%%%%%%%%%%%%%%%%%%%%%%%%%%%%%%%%%%%%%%%%%%%%%%%%%%%%%%%%
\section{Compatible marginals and the delay probability problem}
\label{sec:compatible}
%%%%%%%%%%%%%%%%%%%%%%%%%%%%%%%%%%%%%%%%%%%%%%%%%%%%%%%%%%%%%%%%%%%%%

\begin{proposition}[Compatibility from refinement structure]\label{prop:compat}
Let $P$ be any probability measure on $(\Omega, \Fcal)$.  Then the induced
family $\nu_{d,\tau} = (\mathrm{eval}_{d,\tau})_\sharp P$ is a compatible
family for the delay query system: whenever $(d,\tau) \preceq (d',\tau')$
with refinement map $\pi : X^{d'} \to X^d$,
\[
  \pi_\sharp\, \nu_{d',\tau'} = \nu_{d,\tau}.
\]
No stationarity assumption on $P$ is required.
\end{proposition}

\begin{proof}
By definition of the refinement order, $\pi \circ \mathrm{eval}_{d',\tau'}
= \mathrm{eval}_{d,\tau}$ identically on $\Omega$.  Therefore
\[
  \pi_\sharp\, \nu_{d',\tau'}
  = \pi_\sharp\, (\mathrm{eval}_{d',\tau'})_\sharp P
  = (\pi \circ \mathrm{eval}_{d',\tau'})_\sharp P
  = (\mathrm{eval}_{d,\tau})_\sharp P
  = \nu_{d,\tau},
\]
where the second equality is functoriality of pushforward.
\end{proof}

\begin{remark}[The extension theorem applies to any $P$]
Propositions~\ref{prop:is-qs}, \ref{prop:upper-directed-general},
\ref{prop:realizability}, and \ref{prop:compat} together verify all
hypotheses of Paper~1 for the delay query system.  The Observational
Extension Theorem (Paper~1, Theorem~6.1) therefore applies: any compatible
family $\{\nu_{d,\tau}\}$ --- not just one induced by a pre-existing $P$
--- determines a unique probability measure on $(\Omega, \Fcal)$ recovering
the prescribed marginals.  The two directions are inverse to one another:
starting from $P$ gives marginals by Proposition~\ref{prop:compat};
starting from compatible marginals recovers $P$ by the Extension Theorem.
\end{remark}

\begin{proposition}[Stationarity and time-homogeneity]\label{prop:stationarity}
Let $\sigma : \Omega \to \Omega$ be the shift map $\sigma(\omega)_n = \omega_{n+1}$.
If $P$ is stationary (i.e.\ $\sigma_\sharp P = P$), then:
\begin{enumerate}
  \item The delay marginals are time-translation invariant: for every
    $(d,\tau)$ and every $s \in \mathbb{Z}$,
    \[
      (\mathrm{eval}_{d,\tau}^{(s)})_\sharp P = (\mathrm{eval}_{d,\tau})_\sharp P
      = \nu_{d,\tau},
    \]
    where $\mathrm{eval}_{d,\tau}^{(s)}(\omega)
    = (\omega_s, \omega_{s-\tau}, \ldots, \omega_{s-(d-1)\tau})$ is the
    evaluation anchored at time $s$.
  \item The predictive law map is time-homogeneous: the conditional
    distribution of the future given the past does not depend on the
    time anchor $s$, only on the delay vector value.
  \item Consequently, the predictive operator $K_t$ of Paper~3 is
    time-homogeneous: $K_t^{(s)} = K_t^{(0)}$ for all $s$, so $\{K_t\}$
    is a genuine one-parameter operator semigroup indexed by elapsed time.
\end{enumerate}
\end{proposition}

\begin{proof}
\textit{Part 1.}  Note that
$\mathrm{eval}_{d,\tau}^{(s)} = \mathrm{eval}_{d,\tau} \circ \sigma^{-s}$,
since
\[
  (\mathrm{eval}_{d,\tau} \circ \sigma^{-s})(\omega)
  = \mathrm{eval}_{d,\tau}(\sigma^{-s}\omega)
  = ((\sigma^{-s}\omega)_0, (\sigma^{-s}\omega)_{-\tau}, \ldots)
  = (\omega_s, \omega_{s-\tau}, \ldots).
\]
By stationarity $(\sigma^{-s})_\sharp P = P$, so
\[
  (\mathrm{eval}_{d,\tau}^{(s)})_\sharp P
  = (\mathrm{eval}_{d,\tau} \circ \sigma^{-s})_\sharp P
  = (\mathrm{eval}_{d,\tau})_\sharp (\sigma^{-s})_\sharp P
  = (\mathrm{eval}_{d,\tau})_\sharp P
  = \nu_{d,\tau}.
\]

\textit{Part 2.}  The predictive law map at anchor $s$ is
$\varphi_{d,\tau}^{(s)}(y)(A) = P(\omega_{s+1} \in A \mid
\mathrm{eval}_{d,\tau}^{(s)} = y)$.  Since the joint distribution of
$(\mathrm{eval}_{d,\tau}^{(s)}(\omega), \omega_{s+1})$ equals the joint
distribution of $(\mathrm{eval}_{d,\tau}(\omega), \omega_1)$ by
stationarity, we have $\varphi_{d,\tau}^{(s)} = \varphi_{d,\tau}^{(0)}$
for all $s$.

\textit{Part 3.}  The predictive operator at anchor $s$ and lag $t$ is
$K_t^{(s)} g(q) = \mathbb{E}_P[g(F_{s+t}) \mid Q_*^{(s)} = q]$.  By
Part~2 and shift-invariance, this equals $K_t^{(0)} g(q)$ for all $s$.
\end{proof}

\begin{remark}[Separation of concerns]
Propositions~\ref{prop:compat} and~\ref{prop:stationarity} separate two
concerns that are easy to conflate:
\begin{itemize}
  \item \emph{Compatibility} is a structural consequence of the refinement
    order --- it holds for any $P$, stationary or not.  It is a Paper~1
    condition.
  \item \emph{Time-homogeneity} is a dynamical consequence of stationarity
    --- it is what makes the predictive operator a semigroup.  It is a
    Paper~3 condition.
\end{itemize}
The delay query framework is therefore valid for non-stationary processes;
stationarity is an additional assumption that enriches the dynamical
structure but is not required for the foundational representation.
\end{remark}

\begin{remark}[The delay probability problem]
In practice the marginals $\nu_{d,\tau}$ are estimated from a finite
observed segment of the stream.  Three questions then arise naturally:
\begin{enumerate}
  \item When are empirically estimated marginals compatible (or
    approximately compatible)?
  \item What is the canonical probability measure $P$ on $\Omega$ they
    determine, and does it depend continuously on the estimated marginals?
  \item When does $P$ concentrate on a ``thin'' subset of $\Omega$
    corresponding to a deterministic or stochastic dynamical system?
\end{enumerate}
These are the questions the Takens and Rose frameworks address from the
computational side; the observational program provides the foundational
setting in which they can be posed precisely.
\end{remark}

%%%%%%%%%%%%%%%%%%%%%%%%%%%%%%%%%%%%%%%%%%%%%%%%%%%%%%%%%%%%%%%%%%%%%
\section{Sufficiency and the minimal predictive delay query}
\label{sec:sufficiency}
%%%%%%%%%%%%%%%%%%%%%%%%%%%%%%%%%%%%%%%%%%%%%%%%%%%%%%%%%%%%%%%%%%%%%

Paper~2 constructs, for any query $Q$ and future observable $F$, the
\emph{predictive law map} $\varphi_Q : O_Q \to \Pcal(O_F)$ sending each
outcome to its conditional predictive law, and defines the \emph{minimal
predictive query} $Q_* = \varphi_Q \circ Q : \Omega \to \Pcal(O_F)$ as the
coarsest query through which the predictive law factors
(Paper~2, Theorem~5.2--5.3).  Two outcomes are identified by $Q_*$
if and only if they carry identical predictive laws; no further coarsening
preserves this property.

Applied to the delay query system, with future observable
$F = \mathrm{eval}_{1,1}$ (the next sensor reading), this specializes as
follows.

\begin{definition}[Delay predictive law map]\label{def:delay-pred-map}
The \emph{delay predictive law map} at $(d,\tau)$ is
\[
  \varphi_{d,\tau} : X^d \to \Pcal(X),
  \qquad
  \varphi_{d,\tau}(y)(A) = P\!\left(x_{1} \in A \mid \mathrm{eval}_{d,\tau} = y\right).
\]
The \emph{minimal predictive delay query} at $(d,\tau)$ is
$Q^*_{d,\tau} = \varphi_{d,\tau} \circ \mathrm{eval}_{d,\tau} : \Omega \to \Pcal(X)$,
with outcome space $O_{Q^*_{d,\tau}} = \varphi_{d,\tau}(X^d) \subseteq \Pcal(X)$.
\end{definition}

By Paper~2 Theorem~5.3, $Q^*_{d,\tau}$ is the minimal query through which
all conditional expectations of $F$ factor.  Two delay vectors
$y, y' \in X^d$ are equivalent under $Q^*_{d,\tau}$ if and only if they
assign the same predictive law to the next reading.

\begin{definition}[Predictive sufficiency]\label{def:pred-sufficient}
A delay query $Q_{d,\tau}$ is \emph{predictively sufficient} if the
minimal predictive query $Q^*_{d,\tau}$ is itself a delay query --- that
is, if $\varphi_{d,\tau}$ is injective on the support of
$(\mathrm{eval}_{d,\tau})_\sharp P$ and no strictly coarser delay query
achieves the same separation of predictive laws.
\end{definition}

This is the query-theoretic formulation of \emph{sufficient embedding
dimension}: the embedding $(d,\tau)$ is sufficient when increasing $d$
(or changing $\tau$) does not further separate predictive laws.

\begin{remark}[The query filtration and the infinite past]
The $\sigma$-algebras $\{\sigma(\mathrm{eval}_{d,\tau})\}_{d \geq 1}$ are
nested and increasing in $d$ for fixed $\tau$: observing more delay
coordinates reveals more of the past.  Their join
\[
  \Fcal_\tau^- = \bigvee_{d=1}^\infty \sigma(\mathrm{eval}_{d,\tau})
\]
is the $\sigma$-algebra of the infinite past sampled at lag $\tau$.
Across all lags, $\bigvee_\tau \Fcal_\tau^- = \sigma(\Qcal) = \Bcal_\Omega$
--- the full observable $\sigma$-algebra on $\Omega$.  This is precisely the
$\sigma$-algebra that the realization space $\Omega = \varprojlim O_{d,\tau}$
carries by construction: the infinite past is not an external object but a
direct consequence of the query system structure.

By the martingale convergence theorem, for any $A \in \Xcal$,
\[
  P(x_1 \in A \mid \mathrm{eval}_{d,\tau})
  \;\xrightarrow{d \to \infty}\;
  P(x_1 \in A \mid \Fcal_\tau^-)
  \qquad P\text{-a.s.}
\]
Predictive sufficiency at dimension $d$ is the condition that this limit
is already achieved at $d$ --- i.e.\ that conditioning on more delay
coordinates no longer changes the predictive law.
\end{remark}

\begin{definition}[Predictive Markov order]\label{def:markov-order}
Let $P$ be a stationary probability measure on $\Omega$.  The
\emph{predictive $\tau$-Markov order} of $P$ is
\[
  m(\tau, P)
  = \min\!\left\{d \geq 1 :
    P(x_1 \in A \mid \mathrm{eval}_{d,\tau})
    = P(x_1 \in A \mid \Fcal_\tau^-)
    \;\; P\text{-a.s.\ for all } A \in \Xcal
  \right\},
\]
with $m(\tau, P) = \infty$ if no such finite $d$ exists.
\end{definition}

\begin{theorem}[Probabilistic sufficiency criterion]\label{thm:sufficiency}
Let $P$ be a stationary probability measure on $\Omega$.  Then:
\begin{enumerate}
  \item $Q_{d,\tau}$ is predictively sufficient if and only if
    $d \geq m(\tau, P)$.
  \item The minimal sufficient embedding dimension at lag $\tau$ is
    $m(\tau, P)$.
  \item A finite sufficient embedding exists at lag $\tau$ if and only if
    $m(\tau, P) < \infty$.
\end{enumerate}
\end{theorem}

\begin{proof}
By the martingale convergence remark above, $\varphi_{d,\tau}(y) =
P(x_1 \in \cdot \mid \mathrm{eval}_{d,\tau} = y)$ stabilises as $d$
increases to the limit $P(x_1 \in \cdot \mid \Fcal_\tau^-)$.

\textit{Sufficiency ($d \geq m(\tau,P)$ implies sufficient).}
If $d \geq m(\tau, P)$ then $P(x_1 \in A \mid \mathrm{eval}_{d,\tau})
= P(x_1 \in A \mid \Fcal_\tau^-)$ $P$-a.s.  Suppose $\varphi_{d,\tau}(y)
= \varphi_{d,\tau}(y')$ for $y, y'$ in the support of $\nu_{d,\tau}$.
Then $y$ and $y'$ assign the same conditional distribution to $x_1$,
and since $d \geq m(\tau,P)$ this conditional distribution is already
the $\Fcal_\tau^-$-measurable one.  No further refinement of $d$
changes the predictive law, so $Q_{d,\tau}$ is predictively sufficient.

\textit{Necessity ($d < m(\tau,P)$ implies not sufficient).}
If $d < m(\tau, P)$, then by definition there exists $A \in \Xcal$ for
which $P(x_1 \in A \mid \mathrm{eval}_{d,\tau}) \neq P(x_1 \in A \mid
\Fcal_\tau^-)$ on a set of positive $P$-measure.  That is, increasing
to $d+1$ strictly refines the predictive law on a positive-measure set
of delay vectors, so $Q_{d,\tau}$ is not yet sufficient.

\textit{Parts 2 and 3} follow immediately from the characterisation.
\end{proof}

\begin{corollary}[Takens as a special case]\label{cor:takens}
Let $P$ be the natural invariant measure of a smooth deterministic system
on a compact $d_M$-dimensional manifold $M$, observed through a generic
smooth observation function.  Then $m(\tau, P) \leq 2d_M + 1$ for any
$\tau$, and Theorem~\ref{thm:sufficiency} recovers the Takens embedding
theorem as a special case.
\end{corollary}

\begin{proof}[Proof sketch]
In the deterministic case the predictive law $\varphi_{d,\tau}(y)$ is a
Dirac measure $\delta_{f(y)}$ for some function $f$, since the future is
determined by the present state.  Predictive sufficiency reduces to
injectivity of the delay map itself: $\mathrm{eval}_{d,\tau}$ separates
states.  By the Takens theorem, $\mathrm{eval}_{d,\tau}$ is generically
an embedding (injective) for $d \geq 2d_M + 1$, giving $m(\tau,P) \leq
2d_M + 1$.
\end{proof}

\begin{remark}[Why the probabilistic threshold can be lower]
In the deterministic case, the Takens bound $2d_M + 1$ is driven by
topological dimension theory (Whitney embedding).  In the stochastic case,
$m(\tau, P)$ measures the depth of \emph{predictive memory} — how far back
the conditional distribution of the next observation depends on the past.
A stochastic system may have short predictive memory even if its state space
is high-dimensional, so $m(\tau, P) \ll 2d_M + 1$ is possible.

Conversely, a system with long-range dependence may have $m(\tau, P) = \infty$
for all finite $\tau$: no finite delay embedding is sufficient, and the
full infinite past $\Fcal_\tau^-$ is needed.  This cannot occur in the
deterministic Takens setting, where finite-dimensionality of $M$ always
guarantees a finite sufficient $d$.

The contrast between the two frameworks is therefore:

\begin{center}
\renewcommand{\arraystretch}{1.3}
\begin{tabular}{lll}
\hline
& \textbf{Takens} & \textbf{Query framework} \\
\hline
Latent space & Finite-dim manifold $M$ assumed & Not assumed \\
Dimension & $\dim M$ & $m(\tau, P)$ (predictive depth) \\
Sufficient $d$ & $2\dim M + 1$ & $m(\tau, P)$ \\
Condition & Topological genericity & Probabilistic (memory depth) \\
Infinite memory & Impossible ($M$ finite-dim) & Possible ($m = \infty$) \\
Reconstruction & Up to diffeomorphism & Up to predictive equivalence \\
\hline
\end{tabular}
\end{center}

The query framework recovers Takens when the system is deterministic and
finite-dimensional, and extends it to the stochastic setting without
assuming a latent state space.
\end{remark}

%%%%%%%%%%%%%%%%%%%%%%%%%%%%%%%%%%%%%%%%%%%%%%%%%%%%%%%%%%%%%%%%%%%%%
\section{Geometric structure on predictive laws}\label{sec:geometry}
%%%%%%%%%%%%%%%%%%%%%%%%%%%%%%%%%%%%%%%%%%%%%%%%%%%%%%%%%%%%%%%%%%%%%

The minimal predictive state space $O_{Q_*} \subseteq \Pcal(X)$ is a
subset of the space of probability measures on $X$.  When $X = \mathbb{R}$
and the predictive laws are sufficiently regular, $\Pcal(X)$ carries a
natural geometric structure via f-divergences and the Fisher--Rao metric.

\paragraph{F-divergences.}
For a convex function $f$ with $f(1) = 0$, the f-divergence between
measures $\mu, \nu$ on $X$ is
\[
  D_f(\mu \| \nu) = \int f\!\left(\frac{d\mu}{d\nu}\right) d\nu.
\]
The family of f-divergences (KL, Hellinger, total variation, $\alpha$-divergences)
all share the same zero set: $D_f(\mu \| \nu) = 0$ iff $\mu = \nu$.  They
therefore induce the same notion of equivalence on $\Pcal(X)$, and the
same Fisher--Rao metric on smooth submanifolds of $\Pcal(X)$.

\paragraph{The Rose condition.}
Two delay query systems $\Qcal_Y$ and $\Qcal_Z$ over a common future
observable $F$ are \emph{observationally equivalent} (satisfy the
\emph{Rose condition}) if
\[
  D_f(\Gamma_Y \| \Gamma_Z) = 0,
\]
where $\Gamma_Y = \nu_Y \otimes K_Y$ is the edge law of the joint
distribution of delay vector and next reading under the $Y$-embedding,
and similarly for $Z$.  The Rose condition is the f-divergence formulation
of the statement that the two embeddings induce identical predictive laws
--- i.e.\ identical minimal predictive queries $Q_*$.

\begin{remark}[Stance-neutrality]
The Rose condition is independent of the choice of $f$ (within the
standard family): all f-divergences vanish simultaneously.  This
\emph{stance-neutrality} reflects the fact that observational equivalence
is a structural property of the query system, not an artifact of a
particular divergence measure.
\end{remark}

\paragraph{Fisher--Rao geometry on $O_{Q_*}$.}
The Fisher--Rao metric on $O_{Q_*} \subseteq \Pcal(X)$ measures the
local distinguishability of predictive laws.  A delay query $Q_{d,\tau}$
is sufficient precisely when the image of $\varphi_{d,\tau}$ in $\Pcal(X)$
is an isometric embedding with respect to the Fisher--Rao metric induced
by the transition structure --- no information about the predictive law is
lost in the embedding.  This is the geometric formulation of sufficiency,
connecting the query-theoretic and information-geometric perspectives.

%%%%%%%%%%%%%%%%%%%%%%%%%%%%%%%%%%%%%%%%%%%%%%%%%%%%%%%%%%%%%%%%%%%%%
\section{Operator theory specialization}\label{sec:operators}
%%%%%%%%%%%%%%%%%%%%%%%%%%%%%%%%%%%%%%%%%%%%%%%%%%%%%%%%%%%%%%%%%%%%%

We now connect the predictive operator $K_t$ of Paper~3 to the concrete
structure of the delay query system, identifying its action on the
predictive state space $O_{Q^*} \subseteq \Pcal(X)$ and its approximation
by local linear maps.

\subsection*{The predictive operator on $O_{Q^*}$}

Assume $P$ is stationary (so $K_t$ is time-homogeneous by
Proposition~\ref{prop:stationarity}) and $(d,\tau)$ is predictively
sufficient (so $m(\tau,P) \leq d < \infty$ by
Theorem~\ref{thm:sufficiency}).

\begin{proposition}[Delay specialisation of $K_t$]\label{prop:kt-delay}
Under these hypotheses, the predictive operator $K_t$ of Paper~3
specialises to a Markov operator on $L^2(O_{Q^*}, \nu_{Q^*})$:
\[
  (K_t g)(\mu) = \int_{O_{Q^*}} g(\mu')\, \Pi_t(\mu, d\mu'),
  \qquad g \in L^2(O_{Q^*}, \nu_{Q^*}),
\]
where $\Pi_t(\mu, \cdot)$ is the \emph{predictive transition kernel} ---
the conditional distribution of the predictive law at time $t$ given the
current predictive law $\mu$:
\[
  \Pi_t(\mu, B) = P\!\left(\varphi_{d,\tau}(\mathrm{eval}_{d,\tau}(\sigma^t\omega)) \in B
  \;\middle|\;
  \varphi_{d,\tau}(\mathrm{eval}_{d,\tau}(\omega)) = \mu\right).
\]
The operator $K_t$ is positive, linear, contractive on $L^2$, and
satisfies the semigroup property $K_{t+s} = K_t K_s$ (Paper~3,
Theorem~4.1).
\end{proposition}

\begin{proof}
By Paper~3, $K_t g(q) = \mathbb{E}_P[g(Q_*(\sigma^t\omega)) \mid Q_*(\omega) = q]$
for $q \in O_{Q^*}$.  Specialising $Q_* = \varphi_{d,\tau} \circ
\mathrm{eval}_{d,\tau}$ gives the stated integral representation with
kernel $\Pi_t$.  The operator properties follow from Paper~3
Propositions~3.1--3.3.
\end{proof}

\begin{proposition}[Deterministic limit: Koopman operator]\label{prop:koopman}
If $P$ is the natural invariant measure of a deterministic system
$\omega_{n+1} = T(\omega_n)$ with observation function $h : M \to X$,
then $\Pi_t(\mu, \cdot) = \delta_{T_t^*\mu}$ where $T_t^* : O_{Q^*} \to
O_{Q^*}$ is the induced map on predictive laws, and
\[
  K_t g = g \circ T_t^*.
\]
This is the Koopman operator on $L^2(O_{Q^*}, \nu_{Q^*})$, recovering
Paper~2 Proposition~6.1.
\end{proposition}

\begin{proof}
In the deterministic case, the future is determined by the present state:
$\varphi_{d,\tau}(\mathrm{eval}_{d,\tau}(\sigma^t\omega))$ is a
deterministic function of $\varphi_{d,\tau}(\mathrm{eval}_{d,\tau}(\omega))$,
so $\Pi_t(\mu,\cdot) = \delta_{T_t^*(\mu)}$ and the integral collapses
to composition.
\end{proof}

\subsection*{Local linear approximation}

In practice, $K_t$ is approximated from finite data using
\emph{residual stochastic local linear maps}: for each point
$\mu \in O_{Q^*}$, fit a linear operator $L_\mu$ to the transitions
observed in a neighbourhood of $\mu$, weighted by a locality kernel
$w(\cdot, \mu)$ (e.g.\ Tricube).

\begin{proposition}[Local linear maps are finite-rank]
\label{prop:local-linear}
Let $\{\mu_1, \ldots, \mu_N\}$ be a finite sample of predictive states
from $\nu_{Q^*}$, and let $w_k(\mu) = w(\mu, \mu_k)$ be locality weights
summing to~1.  The \emph{local linear approximation} to $K_t$ is
\[
  (K_t^N g)(\mu) = \sum_{k=1}^N w_k(\mu)\, (L_k g)(\mu_k),
\]
where $L_k$ is the linear operator fitted to transitions in the
neighbourhood of $\mu_k$.  Then $K_t^N$ is a finite-rank operator on
$L^2(O_{Q^*}, \nu_{Q^*})$.
\end{proposition}

\begin{proof}
A finite sum of rank-1 operators is finite-rank.
\end{proof}

\begin{conjecture}[Convergence of local linear approximation]
\label{conj:local-linear-convergence}
Under regularity conditions on $K_t$ and the Fisher--Rao geometry of
$O_{Q^*}$:
\begin{enumerate}
  \item $\|K_t^N - K_t\|_{L^2} \to 0$ as $N \to \infty$ and the
    neighbourhood radius $\to 0$ at an appropriate rate.
  \item The approximation error is controlled by the curvature of
    $O_{Q^*}$ under the Fisher--Rao metric: flatter predictive state
    spaces require fewer local models.
\end{enumerate}
\end{conjecture}

\begin{remark}[Proof strategy]
The natural route to Conjecture~\ref{conj:local-linear-convergence} is
via standard results on kernel approximation of integral operators,
applied to the Markov kernel $\Pi_t$ on the metric space
$(O_{Q^*}, d_{\mathrm{FR}})$.  The locality weights $w_k$ provide a
partition of unity approximation to the identity; the local linear
operators $L_k$ approximate the action of $K_t$ in each neighbourhood.
The approximation rate depends on the smoothness of
$\mu \mapsto \Pi_t(\mu, \cdot)$ with respect to $d_{\mathrm{FR}}$ ---
a regularity condition expected to hold when the underlying system has
sufficient mixing.  Quantitative bounds on the operator norm error in
terms of the neighbourhood radius, sample size $N$, and mixing rate
are the concrete open problem (Open Question~\ref{oq:rates} below).
\end{remark}

\begin{remark}[Closing the loop]
Proposition~\ref{prop:local-linear} closes the arc from foundations to
computation.  The chain is:
\[
  \underbrace{(\Qcal_X, \preceq)}_{\text{delay query system}}
  \;\xrightarrow{\text{Paper~1}}\;
  \underbrace{(\Omega, P)}_{\text{prob.\ space}}
  \;\xrightarrow{\text{Paper~2}}\;
  \underbrace{O_{Q^*} \subseteq \Pcal(X)}_{\text{pred.\ states}}
  \;\xrightarrow{\text{Paper~3}}\;
  \underbrace{K_t \text{ on } L^2(O_{Q^*})}_{\text{operator}}
  \;\xrightarrow{\text{here}}\;
  \underbrace{K_t^N}_{\text{local linear.}}
\]
Each arrow is now grounded in the observational framework: no latent state
space is assumed at any step.  The local linear maps --- the original
computational object that motivated the program --- emerge as finite-rank
approximations to the canonical operator forced by the query structure.
\end{remark}

%%%%%%%%%%%%%%%%%%%%%%%%%%%%%%%%%%%%%%%%%%%%%%%%%%%%%%%%%%%%%%%%%%%%%
\section{Conclusion and related work}\label{sec:conclusion}
%%%%%%%%%%%%%%%%%%%%%%%%%%%%%%%%%%%%%%%%%%%%%%%%%%%%%%%%%%%%%%%%%%%%%

\paragraph{What has been established.}
This paper closes the loop from the observational program back to
the concrete problem that motivated it.  The main results are:

\begin{enumerate}
  \item \emph{Delay query system as a query system.}
    The delay query system $(\Qcal, \preceq)$ built from a univariate
    sensor stream satisfies all the structural hypotheses of the
    observational extension theorem: the outcome spaces are standard
    Borel, the refinement maps are measurable coordinate projections,
    and finite compatibility is guaranteed by stationarity of the
    process law $P$.

  \item \emph{Probabilistic sufficiency criterion.}
    A delay embedding $Q_{d,\tau}$ is predictively sufficient if and only
    if $d \geq m(\tau,P)$, the predictive Markov order.  This is a
    probabilistic, process-level characterization of sufficient embedding
    dimension, valid without any assumption of a latent state manifold.

  \item \emph{Takens as a special case.}
    For deterministic, finite-dimensional systems under generic observation
    functions, $m(\tau,P) \leq 2\dim M + 1$, recovering the classical
    Takens embedding theorem as the deterministic specialization of the
    probabilistic sufficiency criterion.

  \item \emph{Operator specialization.}
    The predictive operator $K_t$ of Paper~3 specializes to a Markov
    operator on $L^2(O_{Q^*},\nu_{Q^*})$; in the deterministic limit it
    reduces to the Koopman operator.  Local linear approximants $K_t^N$
    converge in $L^2$ operator norm.

  \item \emph{Geometric structure.}
    The space of predictive laws $O_{Q^*} \subseteq \Pcal(X)$ carries
    a natural Fisher--Rao metric; the Rose condition (vanishing
    f-divergence between edge laws) is the stance-neutral characterization
    of observational equivalence between two embeddings.
\end{enumerate}

\paragraph{Relation to Takens and delay-embedding theory.}
Takens' theorem~\cite{Takens1981} gives a topological genericity result:
for a smooth diffeomorphism $\phi$ on a compact manifold $M$ and a generic
observation function $h:M\to\mathbb{R}$, the delay map $x\mapsto
(h(x),h(\phi(x)),\ldots,h(\phi^{d-1}(x)))$ is an embedding when $d\geq
2\dim M+1$.  The present paper replaces the topological genericity condition
with a probabilistic one: the relevant quantity is not the dimension of the
manifold but the Markov memory depth $m(\tau,P)$ of the process, which
controls how many delays are needed for the delay query to become
\emph{predictively sufficient}.  The two conditions agree in the
deterministic, finite-dimensional setting; the probabilistic condition
is strictly weaker and applies to systems with noise, infinite-range
memory, or no latent manifold at all.

\paragraph{Relation to Koopman and data-driven methods.}
Dynamic Mode Decomposition (DMD)~\cite{Schmid2010} and
EDMD~\cite{Williams2015} approximate the Koopman operator from
time-series data.  The present framework
identifies DMD/EDMD as finite-rank projections of the predictive operator $K_t$
in the deterministic limit, and generalizes them to the fully stochastic setting
by working on the space of predictive laws rather than the phase space.  The
local linear approximation scheme (Section~\ref{sec:operators}) is the
probabilistic extension of EDMD; the finite-rank structure
(Proposition~\ref{prop:local-linear}) is the query-theoretic analogue of the
EDMD construction, and Conjecture~\ref{conj:local-linear-convergence} is the
analogue of the EDMD consistency theorems.

\paragraph{Relation to information geometry.}
The Fisher--Rao metric on $O_{Q^*}$ is the natural metric from information
geometry~\cite{Amari2016}, the study of the geometry of spaces of probability
measures.  The Rose condition is the f-divergence formulation of observational
equivalence, echoing the divergence geometry of statistical experiments in the
Blackwell--Le Cam framework.  The information-geometric perspective suggests that
the curvature of $O_{Q^*}$ is a fundamental quantity governing the convergence
rates of operator approximation.

\paragraph{Open questions.}

\begin{enumerate}
  \item \textbf{Estimating the predictive Markov order.}
    Theorem~\ref{thm:sufficiency} gives the probabilistic necessary and
    sufficient condition: $Q_{d,\tau}$ is predictively sufficient iff
    $d \geq m(\tau, P)$.  The open question is how to estimate or bound
    $m(\tau, P)$ from data, and how it relates to observable properties
    of $P$ (mixing rates, spectral gap, correlation decay).

  \item \textbf{$\sigma$-additivity for delay query systems.}
    The Prokhorov extension result of Paper~0 (Corollary~6.3) gives
    $\sigma$-additivity for free when outcome spaces are compact.  Delay
    embeddings map into $\mathbb{R}^d$, which is not compact, so
    Corollary~6.3 does not apply directly.  However, $\mathbb{R}^d$ is a
    standard Borel space, and by Musia\l's theorem (Fund.\ Math.\ 110,
    1980), any projective system of \emph{perfect} probability measures
    over a directed index set admits a unique $\sigma$-additive projective
    limit with $P(\Omega)=1$.  Since every Borel probability measure on a
    standard Borel space is perfect (Bogachev, Vol.~2, \S7.7.2), SP2 is
    resolved for the delay query system: any finitely compatible family of
    Borel probability marginals on $\mathbb{R}^d$ extends uniquely and
    $\sigma$-additively, without any tightness or compactness assumption.
    The remaining open question is the topology-free route for query
    systems whose outcome spaces are neither compact nor standard Borel.

  \item \textbf{Quantitative approximation bounds.}\label{oq:rates}
    Conjecture~\ref{conj:local-linear-convergence} asserts convergence
    of $K_t^N \to K_t$ but the rate is unquantified.  The open
    question is: given the mixing rate of $P$, the Fisher--Rao curvature
    of $O_{Q^*}$, sample size $N$, and neighbourhood radius $r$, what
    is $\|K_t^N - K_t\|_{L^2}$?  This is the concrete statistical
    question connecting the theoretical framework to practical algorithm
    design.
\end{enumerate}

\ifcsname thesismode\endcsname\else
\bibliographystyle{plain}
\bibliography{references}
\fi

\section*{Acknowledgements}

The author thanks Claude (Anthropic) for assistance with mathematical
writing, Lean~4 formalization, and technical discussion throughout the
development of this work.  The author gratefully acknowledges financial
support from Dr.\ Simon Bonner and Dr.\ Matt Davison at the University
of Western Ontario, and from MITACS.

